\section{Introduction}
\label{sec:introduction}

Vision-language-action models (VLAs) achieve strong manipulation performance by fine-tuning large language models on robot demonstration data~\cite{brohan2023_rt2,kim2024_openvla}.
The prevailing architectural assumption is that LLM fine-tuning injects action-relevant representational knowledge into the system---that the language backbone learns to transform visual input into representations suited for action prediction.
Yet this assumption has not been tested at the representational level.
Where do the representations that drive action prediction actually originate?
The answer has direct consequences for VLA design: if the LLM backbone is the primary source of action structure, scaling and training the LLM is the critical lever; if instead the pretrained vision encoder already supplies this structure, investment should prioritize visual representation quality.

Recent work has begun to probe this question from complementary angles.
At the behavioral level, \cite{grant2026} demonstrated that VLA policies are disproportionately driven by visual input, with language instructions having limited influence on action selection.
Representational analyses of VLA internals further suggest that language instructions are largely ignored by the action generation pathway~\cite{haon2025}.
However, these studies operate at the behavioral or single-component level---none has performed a systematic \emph{representational decomposition} that quantifies how each architectural component (vision encoder, projector, LLM layers) independently contributes to action-relevant structure.
Without such a decomposition, it remains unclear whether VLA fine-tuning creates new action knowledge or merely reconfigures structure inherited from pretrained components.

We address this gap by decomposing OpenVLA-7B~\cite{kim2024_openvla,karamcheti2024}, a VLA that chains a frozen SigLIP vision encoder~\cite{zhai2023_siglip}, an MLP projector, and a Llama-2 7B language backbone fine-tuned on LIBERO-10~\cite{liu2023_libero}.
We extract representations from the projector output, each of the 32 LLM layers at both distributed image-token and readout-token positions, and two external baselines---CLIP ViT-L/14~\cite{radford2021_clip} and DINOv2 ViT-L/14~\cite{oquab2024_dinov2}---that isolate the effect of pretraining objective (vision-language vs.\ vision-only).
To disentangle training from architecture, we cross two factors in a $2 \times 2$ factorial design: LLM training status (trained vs.\ untrained, with the untrained control reinitializing all 6.74B LLM parameters) and instruction presence (with vs.\ without language instruction).
We quantify action-relevant structure using Representational Similarity Analysis (RSA)~\cite{kriegeskorte2008}, which measures geometric alignment between representations and the action space, and functional MLP probes~\cite{alain2017}, which measure decodability of action vectors from frozen representations.

Three findings emerge from this decomposition.
\emph{First}, the frozen vision-language encoder provides the dominant source of action-relevant geometric structure.
SigLIP+Projector achieves the highest action-space alignment (RSA $\rho = 0.615$); CLIP, an independent VL-pretrained encoder with no VLA training signal, achieves $\rho = 0.557$---both far exceeding vision-only DINOv2 ($\rho = 0.388$).
This 0.169-point gap between CLIP and DINOv2, two ViT-L/14 architectures differing primarily in pretraining objective, isolates VL pretraining as the operative factor.
The LLM never exceeds $\rho = 0.501$ at any layer.
\emph{Second}, the LLM degrades distributed geometric structure ($\rho$: $0.615 \to 0.501$) but concentrates a functionally superior representation at the readout token (probe $R^2 = 0.731$ vs.\ $0.706$ for the projector, $p = 0.001$), characterizing the LLM as an information \emph{router} rather than a knowledge source.
\emph{Third}, the geometric structure is invariant to LLM training ($\Delta\rho = 0.021$, $p = 0.20$, n.s.), while functional concentration at the readout requires it ($\Delta R^2 = 0.045$, $p < 0.0001$).
Training teaches routing, not representation creation.

These findings make three contributions:
\begin{itemize}[nosep,leftmargin=*]
  \item We provide the first representational decomposition of a VLA, separately quantifying the action-relevant structure contributed by the vision encoder, projector, and LLM layers using both geometric (RSA) and functional (probe $R^2$) metrics.
  \item We demonstrate that VL-pretrained vision encoders---not LLM fine-tuning---are the primary source of action-relevant geometric structure, with two independent VL encoders (SigLIP, CLIP) both substantially exceeding vision-only DINOv2.
  \item We characterize the LLM as an information router that concentrates pre-existing vision-encoder geometry at the readout token for action decoding, with training teaching the routing policy rather than creating new representations.
\end{itemize}
These results suggest that vision encoder pretraining quality---specifically, vision-language alignment---is the primary lever for action representations in current VLA architectures, and that the LLM backbone's role is information routing rather than knowledge injection (\S\ref{sec:discussion}).
